\chapter{Thầy Cũng Quên, Huống Gì Trò}
\markboth{Thầy Cũng Quên, Huống Gì Trò}{Thầy Cũng Quên, Huống Gì Trò}

\begin{truyen}{Thầy Vào Trễ Mười Lăm Phút}{The Teacher Gets Caught}
\chuhoa{C}{ả lớp đang rì rầm thì thầy bước vào muộn. Một bạn đứng lên, lễ phép nhưng sắc lắm: ``Thầy dặn tụi em không được trễ một phút. Sao hôm nay thầy vào lớp trễ mười lăm phút ạ?''}

\teacherpair{Thầy cười nửa miệng: ``Thầy vừa họp gấp với nhà trường. Nhưng em hỏi vậy là đúng. Nội quy không tự dưng mất hiệu lực chỉ vì người vi phạm là thầy.''}{The teacher replies: ``I had an urgent meeting with the school. But your question is fair. Rules do not lose force just because the person breaking them is the teacher.''}

Không khí trong lớp im hẳn. Bạn ấy tưởng thầy sẽ gắt. Nhưng thầy nói tiếp:

\teacherpair{Tuy nhiên, phân biệt cho rõ: có cái trễ vì vô ý thức, có cái trễ vì trách nhiệm khác chen vào. Thầy nợ lớp lời xin lỗi vì đến muộn. Còn các em nợ thầy một điều khác: đừng dùng lỗi của thầy làm cớ hợp thức hóa lỗi của mình.}{``Still, we must distinguish between lateness caused by carelessness and lateness caused by another duty. I owe the class an apology for being late. You owe me something else: do not use my mistake to legalize yours.''}

Thầy đặt tập giáo án xuống bàn, nói chậm hơn thường ngày: ``Người lớn nhiều khi cũng sống trong một cái lịch không phải do mình tự chọn hết. Họp, báo cáo, điện thoại từ phụ huynh, việc trường chen vào nhau. Nhưng bận không phải là lá chắn để khỏi xin lỗi. Bận chỉ là lời giải thích, còn lời xin lỗi vẫn là bổn phận.''

\textit{(The class catches the teacher in inconsistency. He admits the fault but keeps the moral distinction intact.)}
\end{truyen}

\begin{baihoc}
Người lớn biết nhận lỗi thì đáng tôn trọng hơn. Nhưng học sinh cũng không được biến lỗi của người lớn thành giấy phép cho mình sống dễ dãi. Công bằng không phải là xóa chuẩn; công bằng là giữ cùng một nguyên tắc cho cả hai phía.
\textit{(Adults are more respectable when they admit mistakes. But students should not turn adult mistakes into permission slips for their own looseness. Fairness does not erase standards; it applies one principle to both sides.)}
\end{baihoc}

\begin{ghinhoanh}
\textbf{Short English to remember:}

Fairness includes teachers too.

One mistake does not erase all standards.
\end{ghinhoanh}

\begin{tuhocnhanh}
1. Khi người lớn sai, em muốn họ xin lỗi hay muốn lấy đó làm cớ cho mình? \textit{(When adults are wrong, do you want an apology or an excuse for yourself?)}

2. Em có đủ công bằng để nhìn hai lỗi khác bản chất không? \textit{(Are you fair enough to see that two mistakes can differ in nature?)}
\end{tuhocnhanh}

\ngancach